% Section 8: Relation to the literature (compressed from archive/sec8.tex; sources: research/G1/RH/spectral/NOVELTY_Q.md,
% NOVELTY_TODA.md, NOVELTY_ZETA.md). Every work added in the Rule L pass of 2026-10-23 was read only through a summarising fetch
% tool (WebFetch); none is quoted in a CITATIONS.md, hence "according to" and "we have not found" throughout. Rakhmanov,
% DragnevSaff are quoted with hypotheses in research/G2/CITATIONS.md, entries (R96), (DS97); BeckermannKuijlaars2001 has not been
% read. MathSciNet and zbMATH were not available to the searches.
\section{Relation to the literature}\label{sec:literature}

\emph{Spectral data and Jacobi parameters.} Orthogonal polynomials whose moments are spectral zeta values at positive integers are classical for
Bessel zeros: the Lommel polynomials are orthogonal on the reciprocal zeros \cite{VA} (Proposition~\ref{prop:bessel}), and an analogue for the
zeros of $J_\nu'$ is given in \cite{ChungLeePark}. The Tricomi--Carlitz polynomials, orthogonal on $\pm(k+\alpha)^{-1/2}$ \cite{WuLinXuZhao}, have
explicit recurrence coefficients in which, from the recurrence as stated there, the offset enters only through the index $n+\alpha$, as in
Theorem~\ref{thm:C} (at $p=1$ in $x^2$, outside our range). Positive Jacobi matrices with compact inverse are studied in
\cite{StovicekSwiderski} through a regularised characteristic function; according to our reading, asymptotics of the Jacobi parameters are
not treated there. We know of no expansion of recurrence coefficients in spectral zeta values at non-positive arguments, as in
Theorems~\ref{thm:A}--\ref{thm:C}. For continuous weights, singularities of the weight produce identifiable power-law or oscillating terms in
the recurrence coefficients \cite{Magnus,FoulquieMoreno}; the pole terms of Theorem~\ref{thm:D} and Conjecture~\ref{conj:W} may be viewed as an
instance of this principle at the accumulation point of a discrete spectrum, with the poles of $Z(-w)$ as singularity data. Subleading
corrections to Lanczos coefficients have recently been derived from endpoint data of continuous densities, including a hard-edge $n^{-2}$
term \cite{MuruganVanZylWatanabe}, which we have not compared with $f_2$.

\emph{Constrained equilibrium.} Existence, uniqueness and structure of equilibrium problems with an upper constraint are classical for
finite constraints or compact intervals \cite{Rakhmanov,DragnevSaff}. For constraints of finite mass with compact support and no external
field, the structure of Proposition~\ref{prop:Q} appears in the analysis of conjugate-gradient convergence \cite{BeckermannKuijlaars2001}:
according to \cite[Lemma~1.5]{BeckermannHelart}, which follows \cite[Lemma~3.1(a)]{BeckermannKuijlaars2001}, a monotonicity criterion gives a
single band and an endpoint equation that becomes~\eqref{eq:Q} after a reflection, and according to \cite{BeckermannGuettel,BeckermannKuijlaars2002}
the proof of \cite[Theorem~2.1]{BeckermannKuijlaars2001} gives a constrained analogue of the Buyarov--Rakhmanov formula \cite{BuyarovRakhmanov}.
In our setting that superposition is $\int_a^Y|F'(b)|\,\omega_{[0,b]}\diff b$, which we have checked numerically, not proved, to agree
with~\eqref{eq:nu}. What is new here, as far as we know, is the constraint of infinite mass at a hard edge (with $Y=\infty$ allowed), where the
finite-mass theory does not apply; the single-integral form~\eqref{eq:nu}; and the edge constant. Corollary~\ref{cor:Q1} is an elementary
evaluation for the lattice problem of \cite[Theorem~1.3]{PaperI}. We have not been able to consult \cite{BeckermannKuijlaars2001} itself. We have
not found Proposition~\ref{prop:M12}; its termwise Mellin asymptotics is a standard technique.

\emph{Shifted index, reflection, Toda.} For varying weights the recurrence coefficients have expansions in even powers of $1/n$ after a
half-index shift, by a reindexing symmetry of the string equations \cite{BleherIts}, by the structure of the local parametrices
\cite{KuijlaarsTibboel}, and as generating functions of maps \cite{EMP}; for Jacobi-type weights see \cite{KMVV}, and an exact pairing of
$n+\alpha/2$ with $\alpha^2$ occurs for semiclassical Laguerre weights \cite{FilipukVanAsscheZhang}. Our reflection uses the grading of
\cite[Lemma~4.1]{KuijlaarsTibboel}. We have not found the decoupling of the Szeg\H{o} exponent from the Bessel order, the identity~\eqref{eq:FI}
or its proof from a second frame at a non-integer index. The Toda hierarchy for tilts is classical \cite{VanAsscheToda}; re-expanding it in
$1/n$ to generate expansion coefficients is the method of \cite[Section~4]{EMP} for varying weights; the mechanism of
Theorem~\ref{thm:indexshift} is the genus-zero one of map enumeration \cite{Ercolani,EMP,Bouttier}, with endpoints as Riemann invariants of the
dispersionless hierarchy \cite{KonopelchenkoMAM}. The new ingredients there are the setting and the zeta-valued times, not the mechanism.

\emph{Hurwitz zeta values.} Hurwitz's formula is quoted from \cite[Theorem~1.1]{DixitKumar}, where it is attributed to Hurwitz and quoted from
Titchmarsh's monograph, which we have not consulted. We have not found the index transfer of Section~\ref{sec:hurwitz} or
Theorems~\ref{thm:C}--\ref{thm:D}; the closest exact instances are the Lommel and Tricomi--Carlitz recurrences, in which an offset of the nodes
enters only through the index.
